\documentclass{article}
\usepackage[utf8]{inputenc}
\usepackage[T1]{fontenc}
\usepackage{geometry}
\usepackage{graphicx}
\usepackage{float}
\usepackage{hyperref}
\usepackage[english,indonesian]{babel}

\title{Optimasi Pemrosesan Template LaTeX Menggunakan Pendekatan Berbasis AI}
\author{Academic Support Agent}
\date{\today}

\begin{document}
\maketitle

\begin{abstract}

\end{abstract}

\chapter{PENDAHULUAN}


\chapter{TINJAUAN PUSTAKA}


% Original content (commented out):
% \documentclass[a4paper,onepage]{article}
\usepackage{polyglossia}
\setmainlanguage{spanish}
\usepackage{hyperref}
\urlstyle{same}
\usepackage{fullpage}
\usepackage{amssymb}
\usepackage{xspace}
\usepackage{graphicx}
\usepackage{longtable}
\usepackage[dvipsnames]{xcolor}
\definecolor{darkred}{rgb}{0.55, 0.0, 0.0}
\usepackage{savetrees}
\usepackage{lipsum}
\newcommand{\prog}[1]{\texttt{#1}\xspace}

\title{Título provisional}
\author{autor@micorreo.es}
\date{Borrador 1, \today}

\begin{document}

\maketitle

\tableofcontents

\newpage

\section{Descripción}

Plantilla básica para español que usa polyglossia (compila con xelatex) para documentos de clase artículo sin márgenes de encuadernación y con aprovechamiento de espacio para ahorrar papel.

\section{Primera cosa}

\lipsum

\subsection{Primera cosa: componentes}

\lipsum

\subsubsection{Despiece}

\lipsum

\subsubsection{Mejoras}

\lipsum

\subsection{Primera cosa: conclusión}

\lipsum

\section{Segunda cosa}

\lipsum

\end{document}


\end{document}